%% -----------------------------------------------------------------------------------------------
\section{Commanded versus executed motion (Check 4)}\label{sec:s-check4}

\paragraph{Definition} Check 4 (item 4, ``motion follows command'', of the audit in \path{src/hci/trajectory_audit.py}) was declared before the C1 run. For each move step $j$, let $\mathbf{c}_j$ be the commanded horizontal TCP increment and $\mathbf{a}_j$ the executed horizontal TCP displacement between consecutive logged poses. A move passes when every step has $\cos(\mathbf{a}_j,\mathbf{c}_j)\ge0.9$ and $\|\mathbf{a}_j\|\ge0.2$\,mm. The length condition is on the executed step itself, not on its component along the command.

\paragraph{Results} V2 failed in 9 of 41 (seed 1) and 5 of 39 (seed 2) moves; the official FAIL is kept. Three computations (the audit CSV, re-running the audit function, an independent recomputation) agreed on all 715 acted branches of all arms. In every failing move of every arm (110 moves), only the first step failed; in the 14 V2 failures only the direction condition failed (first-step cosine 0.63--0.88, smallest first step 0.47\,mm), and the later steps of the same moves had cosine 0.934--0.999. A post hoc check of steps 2 onwards passed in all 715 moves. The development harm 51157 passed Check 4.

\paragraph{Associations, not causes} The first-step failures had no robot--target or robot--table impulse in the contact list. The failing V2 moves had a lower cosine between the first command and the TCP motion in the step before the trigger (median $-0.57$ against $0.66$ for passing moves) and a lower first-step length (median 1.08 against 1.65\,mm). This is consistent with the arm still carrying the policy's previous motion at the trigger. Controller internals (inverse-kinematics targets, drive targets, sub-step states) were not logged during C1, so the cause of the first-step deviation is not established.
